\chapter{Introduction}

In 2017, a group of researchers working at Google published a paper introducing a technology that revolutionized society.
In \textit{Attention is All You Need}, they presented the Transformer architecture, today's core of large language models (LLMs) and generative AI.

This architecture was built as a text sequence-to-sequence model for machine translation, but today it finds applications in various fields, such as natural language processing, computer vision, and speech recognition.

The first transformer-based model open for public use was ChatGPT by OpenAI.
It made its debut in November 2022 and by December 2024, it reported 300 million global weekly users.
Now there are many models competing for this market, such as Claude, from Anthropic, and Gemini from Google.

These models can all be categorized as generative AI tools.
This technology is arguably the most rapidly adopted in history.
In the second half of 2025, roughly one in six people worldwide was using generative AI tools \citep{microsoft2026aidiffusion}.
LLMs are not only being used as chatbots for answering simple questions, but are also shaping the future of industry and society.
For example, \citet{liang2024mappingincreasingusellms} found that more than 17\% of the papers in computer science published between January 2020 and February 2024 on the \textit{arXiv} and \textit{bioRxiv}, and in \textit{Nature} portfolio journals had sentences heavily edited by ChatGPT. 
Naturally, the question of their computational power arises.
Are these models genuinely capable of performing the tasks for which they are being employed?
There had been different attempts to formally answer this question.
In this thesis, we try to contribute to the community efforts.
We approach the problem from the perspective of formal languages.
We aim to understand the limitations that transformers inherit from their own architecture, which had been optimized for parallelism and learnability, and not necessarily for solving problems in the sense of computability and complexity theory.


\bigskip

First, we need to define what an LLM is.
We consider LLMs simply as instantiations of the transformer architecture.
Thus, we need to understand what transformers are.

A transformer is a machine consisting of a tape from which it reads the input and onto which it appends the output character by character in an autoregressive manner.\footnote{An autoregressive machine computes future values based on the previously computed, i.e., it uses its previous outputs as input for the next step.}
It reads the content of the tape, then computes the next character, appends it at the end, and repeats.
\citet{ATTN-all-you-need} presented a concrete way to do this.

First, the input text written by the user is split into tokens.
These tokens are not necessarily single characters or words of human language, they are sequences of characters previously defined.
The tokenizing process is not of interest since it is just a way of splitting the human text into something the transformers can understand; thus, in this thesis (and in all papers in the area) it is assumed that the user provides input in tokens, i.e., the transformers receive tokens on their input tape.
The set of tokens can be interpreted as the alphabet or vocabulary of the transformers.

Then, the first real component of a transformer is the mechanism that maps the tokens of the tape into vectors of numbers.
Each occupied position of the tape is mapped to a vector.
Each vector is computed depending on its position in the tape and the token at that position.
Thus, this mapping from tape positions to vectors can be decomposed into two components: the position encoding and the token embedding.
The dimension of the vectors is the dimension of the transformer.

Then, these vectors pass through a fixed sequence of layers consisting of an attention mechanism and a feed-forward network.
These mechanisms encode information about how vectors relate to one another, i.e., how a token at a given position relates to a token at another position.
The output of each layer is added to the original vectors and passed to the next layer.

In the vector corresponding to the last token of the tape, after passing through all the layers, a probability distribution over the next token is encoded.
This distribution is extracted and used to pick the new token to append to the tape.
The token is selected, and the process is repeated.
In this way, LLMs like ChatGPT generate the text they return in response to user's prompts

\medskip

The performance of LLMs improves when we allow them to generate more intermediate tokens.
Intuitively, it is like giving them a scratch pad.
This process is known as Chain of Thought, and it is known that increasing the number of intermediate tokens improves the quality of the answers \citep{wei2022chain}.
Notably, using heuristic phrases in prompts, such as ``let's think step by step'', improves the answers without the need to retrain the model \citep{kojima2022large}.
Also, providing the first steps of the reasoning in the prompt, even if they are wrong, gives better results \citep{wang-etal-2023-towards,cot-prompt}.
This is why modern models perform this reasoning in the background before responding to the user.

\bigskip

In this thesis, we analyze transformers as language recognizers.
We define them as machines that receive an input string and determine whether it has a certain property, generating the token $\qaccept$ or $\qreject$, respectively.
When working with Turing machines as language recognizers, it is standard to study the languages recognizable when the machines are limited to a certain time or space budget as a function of the input size.
When working with transformers, there are different resources that are of interest to bound.
For example, \citet{merrill2026little} studied families of transformers whose number of layers grows with the input size.
\citet{merrill2026little} prove that even with a logarithmic growth in the number of layers, there is a significant improvement in the languages that transformers can recognize.

In this work, we extend the results of \citet{toyota}.
They studied the number of steps of Chain of Thought ($\CoT$) as the limited resource.
In their work, they give a formal argument for why transformers perform better with more steps of $\CoT$.
They prove that with $T(n)$ steps of $\CoT$, a transformer can recognize languages recognizable by Boolean circuits of size $T(n)$.
They also prove that a logarithmic number of steps of $\CoT$ does not expand the set of languages recognizable by transformers without $\CoT$.

\citet{toyota} give bounds on the computing power of transformers and define classes of languages depending on how many steps of $\CoT$ are required to recognize them.
They do not study the properties of those classes.
In this thesis, as a preliminary step toward our main objective, we prove that the classes of \citet{toyota} are closed under Boolean operations.

\medskip

Most of the formal work done to understand the computational power of transformers ignores one of their main characteristics.
\citet{toyota} and \citet{merrill2026little} work with deterministic abstractions, but transformers are intrinsically stochastic.
As was mentioned before, in each step of $\CoT$, transformers generate a probability distribution to decide the next token.
Formal studies tend not to sample this distribution and just define the next token as the token with the largest probability in the distribution.
In practice, the models behave in a nondeterministic way by sampling the distribution to define the next token.
The stochasticity of the models is controlled through a parameter called temperature, and it is well known that a nonzero temperature is sometimes necessary for good performance \citep{demir2025optimal}.

\smallskip

In this thesis, we adapt the abstraction of \citet{toyota} to incorporate the stochastic behavior.
We define the corresponding complexity classes and establish the analogues of the deterministic bounds in the probabilistic setting.
We complete the work by studying the properties of the classes of languages recognizable by probabilistic transformers with different numbers of steps of $\CoT$.
The classes defined in this work are closed under Boolean operations and behave well as probabilistic classes, i.e., regardless of how much error we allow the probabilistic transformers to have, the recognizable languages remain the same.

The main contribution of this work is that we prove relations between the power of probabilistic and deterministic transformers.
These relations provide a solid foundation for understanding when it suffices to study deterministic transformers and when a probabilistic abstraction is warranted.






% \begin{itemize}
%     \item well-suited to parallelism enables transformer models to take full advantage of the power and speed offered by GPUs during both training and inference. This possibility, in turn, unlocked the opportunity to train transformer models on unprecedentedly massive datasets through self-supervised learning.
%     \item chatgpt, claude, gemini, 
%     \item every day users
%         \item number of users?
%     \item use in the industry?
%     \item other uses besides the chatbots?
%     \item no usamos normalizacion porque andie lo usa, es complicado. Si se achican los bitsize empieza a tener menos sentido
% \end{itemize}
